\chapter{Resultados y conclusiones}

\section{Resultados}
Tech-Flow cumple su objetivo operativo: autentica por rol, controla la apertura de caja como interruptor del punto de venta, mantiene inventario con avisos en rojo, registra ventas con garantías y entrega informes de ventas y entradas en .txt, además de una vista de inventario filtrable por categoría que se alimenta del propio catálogo. La suite de 516 pruebas pasa en limpio y las vistas cargan sin errores.

\section{Conclusiones}
\begin{itemize}
\item La persistencia CSV simplifica despliegue y respaldo a costa de integridad referencial explícita: la disciplina de escritura (leer, validar, reescribir) se sostiene por las transacciones pequeñas del dominio.
\item La decisión de una sola ventana evita problemas de gestores de ventanas y simplifica el flujo de navegación; los overlays amplían la utilidad de esa decisión.
\item Los últimos meses de mantenimiento mostraron que la duplicación estaba donde crecía más rápido (tablas, tarjetas, guardado .txt), por eso el refactor centralizó helpers compartidos.
\item El bug del precio de venta fue el fallo más grave: la interfaz y el dominio inconvenían \texttt{precioVenta} y la fórmula \texttt{precioEntrada$\times$1.20$+$13\% IVA}; corregido en \texttt{7e73fed}.
\end{itemize}

\section{Recomendaciones}
\begin{itemize}
\item Cifrar contraseñas y sacar \texttt{data/} de git (objetivo de la próxima iteración).
\item Introducir constantes de esquema por archivo para eliminar el acceso por índice numérico.
\item Añadir pruebas tras probar interacciones GUI con TestFX o una capa de servicio testeable.
\item Revisar la unicidad de \texttt{idCliente} para CI y la posible colisión de CI no numéricos.
\end{itemize}
